\section{Алгебра булевых функций и булевы алгебры. Алгебра Линденбаума-Тарского. Эквивалентные преобразования. Разложение Шеннона.}

\subsection*{Эквивалентные преобразования формул}

Для преобразования формул с сохранением их смысла (реализуемой функции) используются два основных правила.

\paragraph{Правило подстановки (Теорема 1)}
Если две формулы равносильны ($\mathcal{J}_1 = \mathcal{J}_2$), то при замене \textbf{всех} вхождений некоторой переменной $x$ на одну и ту же формулу $\mathcal{S}$ результат останется равносильным:
$$\mathcal{J}_1\{ \mathcal{S} // x \} = \mathcal{J}_2\{ \mathcal{S} // x \}.$$
\textit{Важно:} условие замены \textbf{всех} вхождений критично (иначе тождество может нарушиться).

\paragraph{Правило замены (Теорема 2)}
Если в формуле $\mathcal{J}$ заменить некоторую подформулу $\mathcal{S}_1$ на равносильную ей формулу $\mathcal{S}_2$, то полученная формула будет равносильна исходной:
$$\mathcal{S}_1 = \mathcal{S}_2 \implies \mathcal{J}\{ \mathcal{S}_1 \} = \mathcal{J}\{ \mathcal{S}_2 / \mathcal{S}_1 \}.$$

\subsection*{Алгебра булевых функций}

Операции $\lor, \land, \neg$ можно рассматривать не только на значениях $\{0, 1\}$, но и на самих функциях $f, g \in P_n$. Результатом операции $f \lor g$ будет новая функция, значение которой на каждом наборе равно дизъюнкции значений исходных функций.

\paragraph{Алгебра Линденбаума-Тарского}
Это формализация «алгебры логики» как математической структуры.
\begin{itemize}
    \item Рассмотрим множество всех формул над некоторым базисом.
    \item Введем на нем отношение эквивалентности $\equiv$ (равносильность).
    \item Множество классов эквивалентности формул относительно операций $\lor, \land, \neg$ образует \textbf{булеву алгебру} (см. билет 21).
\end{itemize}
Это позволяет нам преобразовывать сложные логические выражения, используя законы булевой алгебры (дистрибутивность, законы де Моргана и т.д.), как обычные алгебраические уравнения.

\subsection*{Разложение Шеннона (Разложение по переменным)}

Любую булеву функцию можно разложить по любой её переменной $x_i$. Это фундаментальная теорема, позволяющая сводить функцию $n$ переменных к функциям $n-1$ переменной.

\paragraph{Теорема (Разложение Шеннона)}
Для любой функции $f(x_1, \dots, x_n)$ справедливо:
$$f(x_1, \dots, x_n) = \left( x_i \land f(x_1, \dots, 1, \dots, x_n) \right) \lor \left( \neg x_i \land f(x_1, \dots, 0, \dots, x_n) \right).$$

\textit{Пояснение:}
\begin{itemize}
    \item $f(x_1, \dots, 1, \dots, x_n)$ — это функция, полученная подстановкой константы 1 вместо $x_i$.
    \item $f(x_1, \dots, 0, \dots, x_n)$ — это функция, полученная подстановкой константы 0 вместо $x_i$.
\end{itemize}

\paragraph{Значение разложения}
1. Позволяет систематически строить формулы для любых функций (через СДНФ).
2. Является основой для работы алгоритмов упрощения булевых функций и построения бинарных диаграмм решений (BDD).

\subsection*{Примеры эквивалентных преобразований}
\begin{enumerate}
    \item $x \to y = \neg x \lor y$ (выражение импликации через базис $\{\neg, \lor\}$).
    \item $\neg(x \land y) = \neg x \lor \neg y$ (закон де Моргана).
    \item $x \lor (x \land y) = x$ (закон поглощения).
\end{enumerate}